% Chapter 6 -- Rubric: PO6 (society/health/safety), PO7 (environment/sustainability),
% PO8 (ethics), PO9 (individual and teamwork), PO11 (project management and finance)
\chapter{Project Management, Teamwork, and Impact}
\label{ch:management}

\section{Project Management and Finance}
\label{sec:management}

\subsection{Planning and Risk Management}
\label{subsec:planning}

\paragraph{Process.} The project was delivered incrementally, as constraint C-01
required: development ran across two academic semesters alongside concurrent
coursework, so no phase could assume a block of uninterrupted time. Work
proceeded on short feature branches merged into \texttt{main} through pull
requests, with the automated pipeline described in \Cref{subsec:eval-setup} as a
required check before merge.

\paragraph{Phases.} Five phases were planned, ordered by dependency rather than
preference, since each supplies the input to the next:

\begin{enumerate}
  \item \textbf{Requirements and connector feasibility.} Elicitation with SSCL
  and study of the candidate tools' public APIs to establish which metrics were
  obtainable in practice. This had to precede any scoring design, because the
  constraint recorded as C-05 --- that several intended metrics are simply not
  exposed --- determines which dimensions the model can support at all.
  \item \textbf{Ingestion.} The connector abstraction, the snapshot data model,
  the job queue and SSE progress streaming.
  \item \textbf{Scoring.} The eight risk strategies, the shared null-aware
  weighted-scoring math, and the score-breakdown presentation.
  \item \textbf{Surveys.} Question generation and quality gating, anonymous
  dispatch and collection, insight generation, and the scheduled monthly
  distribution cycle.
  \item \textbf{Action logging and consolidation.} The action log, lexical
  search with optional AI-assisted reranking, effectiveness review, and the
  portfolio and tracking refinements. INV-3 found that the intended
  reranking-to-lexical failover is not implemented.
\end{enumerate}

The project ran over one year, across two academic semesters.
\Cref{fig:schedule} shows how the phases were distributed across them. The
unplanned work described in \Cref{subsec:deviation-email} (email provider
migration) and \Cref{subsec:deviation-guidance} (removal of manual survey
guidance) was absorbed within the survey phase.

\begin{figure}[H]
\centering
\begin{ganttchart}[
    hgrid, vgrid,
    x unit=2.2cm, y unit chart=0.75cm, y unit title=0.7cm,
    title label font=\small\bfseries,
    bar label font=\small,
    bar/.append style={fill=gray!55},
    bar height=0.5
  ]{1}{4}
  \gantttitle{Semester 1}{2} \gantttitle{Semester 2}{2} \\
  \gantttitle{First half}{1} \gantttitle{Second half}{1}
  \gantttitle{First half}{1} \gantttitle{Second half}{1} \\
  \ganttbar{1. Requirements and feasibility}{1}{1} \\
  \ganttbar{2. Ingestion}{1}{2} \\
  \ganttbar{3. Scoring}{2}{3} \\
  \ganttbar{4. Surveys}{3}{4} \\
  \ganttbar{5. Action logging and consolidation}{4}{4}
\end{ganttchart}
\caption{Phase schedule across the two semesters}
\label{fig:schedule}
\end{figure}

\paragraph{Risk management.} Risks were identified at the start of the ingestion
phase and reviewed at each phase boundary. \Cref{tab:risks} records each with
the mitigation adopted and the exposure that remains after it --- the residual
column being the one that matters, since a risk register listing only
mitigations claims a completeness no project has.

\begin{longtable}{
  >{\raggedright\arraybackslash}p{0.05\textwidth}
  >{\raggedright\arraybackslash}p{0.22\textwidth}
  >{\raggedright\arraybackslash}p{0.32\textwidth}
  >{\raggedright\arraybackslash}p{0.28\textwidth}
}
\caption{Project risks, mitigations and residual exposure}
\label{tab:risks}\\
\toprule
\textbf{ID} & \textbf{Risk} & \textbf{Mitigation adopted} & \textbf{Residual exposure} \\
\midrule
\endfirsthead
\toprule
\textbf{ID} & \textbf{Risk} & \textbf{Mitigation adopted} & \textbf{Residual exposure} \\
\midrule
\endhead
\bottomrule
\endlastfoot
R1 & Third-party rate limits exhausted by portfolio-wide scheduled
synchronisation (C-03, NFR-16) & Bounded connector concurrency (4 per project),
staggering of projects sharing a credential (10\,min default within a 180\,min
spread), pre-flight rate-limit checks in the GitHub connector, and product
positioning as a trend indicator rather than a real-time monitor & A ceiling on
projects per credential that INV-4 could not quantify for want of a
per-connector call count; a portfolio beyond that size needs additional
credentials \\
\addlinespace
R2 & Dependence on external AI providers for embeddings, reranking and text
generation & Ordinary lexical action search remains independent of AI, and
action embedding is asynchronous & A requested deep search returns HTTP 503
rather than falling back when Pinecone is unavailable; survey question
generation and insight have \textbf{no} non-AI substitute, so an extended Gemini
outage postpones a survey cycle outright (NFR-13, INV-3) \\
\addlinespace
R3 & Metrics assumed available but not obtainable in practice (C-05) & API
feasibility study before scoring design; null-aware renormalised scoring so an
absent metric is excluded rather than penalised (FR-22), rationale in
\Cref{subsec:inv1} & Scores computed from fewer signals are weaker estimates and
the interface does not communicate this (\Cref{subsec:deviation-confidence}) \\
\addlinespace
R4 & Free-tier services withdrawn, throttled or altered mid-project & All
environment-specific configuration externalised (NFR-24) so a provider can be
substituted without code change --- which is precisely what made the
Gmail-to-Brevo migration cheap & A provider change still requires a person to
perform it; no fallback provider is pre-configured for any service \\
\addlinespace
R5 & Loss of team capacity to coursework and examination periods (C-01, C-02) &
Small, independently mergeable units of work; trunk kept deployable; CI as the
check that a returning member had not broken something they had forgotten the
details of & Phases dependent on one member's knowledge resumed more slowly; no
dedicated QA or operations capacity existed to absorb the gap \\
\addlinespace
R6 & Partner data unavailable for validation under the NDA (C-09) &
Investigations designed to be answerable without partner data --- INV-1 argues
from the scoring arithmetic and INV-2 from the schema, neither of which requires
a dataset & Absent-metric behaviour (INV-1) is argued analytically rather than
measured on real snapshots; retrieval \emph{quality} (INV-3) and rate-limit
headroom (INV-4) remain unquantified; evaluation used the team's own
repositories rather than partner data \\
\addlinespace
R7 & Anonymity guarantee eroded by later development & Direct identity storage
is prevented by the response schema; shared-token and non-joinable-recipient
behaviour is explicit in the current application design (INV-2) & The
application-level properties can be weakened by later code changes without
changing \texttt{survey\_response}; statistical inference routes A4--A6
remain \\
\end{longtable}

R7's residual column reflects the design principle behind INV-2: guarantees
should be placed in the schema wherever the schema can express them, because
application-level properties such as the shared token and dispatch behaviour
remain choices that a later change could reverse.

\subsection{Budgeting and Resource Identification}
\label{subsec:budget}

Constraint C-03 permitted no monetary budget, so every external service was used
within its free tier. The project incurred no monetary cost
(\Cref{tab:budget}).

\begin{longtable}{
  >{\raggedright\arraybackslash}p{0.22\textwidth}
  >{\raggedright\arraybackslash}p{0.42\textwidth}
  >{\raggedright\arraybackslash}p{0.14\textwidth}
  >{\raggedright\arraybackslash}p{0.06\textwidth}
}
\caption{Resources consumed and their cost}
\label{tab:budget}\\
\toprule
\textbf{Resource} & \textbf{Purpose} & \textbf{Tier} & \textbf{Cost} \\
\midrule
\endfirsthead
\toprule
\textbf{Resource} & \textbf{Purpose} & \textbf{Tier} & \textbf{Cost} \\
\midrule
\endhead
\bottomrule
\endlastfoot
Supabase (managed PostgreSQL) & Primary datastore: snapshots, scores, surveys,
actions, embeddings & Free tier & \$0 \\
\addlinespace
Managed Redis & BullMQ backing store, sessions, and short-lived
reset/invite/verification tokens & Free tier & \$0 \\
\addlinespace
Render & API and always-on worker hosting & Free tier & \$0 \\
\addlinespace
Vercel & Static frontend hosting & Free tier & \$0 \\
\addlinespace
Google Gemini API & Embeddings (\texttt{gemini-embedding-001}, 768 dims) for
action search; question generation and survey insight
(\texttt{gemini-2.5-flash}) & Free tier & \$0 \\
\addlinespace
Pinecone & Reranking of action-search candidates in explicit deep mode & Free
tier & \$0 \\
\addlinespace
Brevo & Transactional email (verification, reset, invitation, survey) & Free
tier & \$0 \\
\addlinespace
GitHub, Jira, SonarCloud & Metric sources & Team's own accounts & \$0 \\
\midrule
\textbf{Total} & & & \textbf{\$0} \\
\end{longtable}

\subsection{Economic Analysis and Financial Prospecting}
\label{subsec:economics}

\paragraph{Operating cost and how it scales.} Operating cost decomposes into
three usage-proportional components and one fixed component, and the shape of
that decomposition is the economically interesting finding.

Per tracked project per month the system consumes: one scheduled synchronisation
per configured slot per day, each producing one snapshot and one scoring
pass --- both of which are compute and storage against the database rather than
calls against a paid API; one survey cycle, costing exactly two Gemini
generation calls (one to generate questions, one to analyse responses); and one
embedding call per logged action, cached by \texttt{content\_hash} so that an
unedited action is never re-embedded. The fixed component is the always-on
worker, which must run continuously regardless of portfolio size.

\textbf{The qualitative conclusion follows without needing the exact numbers:
per-project marginal cost is very low --- two AI generation calls per month plus
database storage --- while the fixed worker cost dominates at small portfolio
sizes.} That cost shape argues strongly for a single multi-tenant hosted
offering over a per-organisation deployment, since the dominant cost is
amortised across all tenants in the first case and replicated per customer in
the second.

At the current scale every one of these components fits within a free tier, so
the operating cost is nil. As a portfolio grows, the always-on worker hosting
and the database are the first components that would require paid plans, since
they are the ones whose free-tier limits are tied to continuous running and
stored volume rather than to occasional calls.

\paragraph{Value proposition.} The benefit to an organisation of SSCL's kind is
the recurring manual effort removed: a manager reconciling four to five tool
dashboards per project on a fixed cadence, which \Cref{ch:problem} identifies as
the originating problem. That effort grows with every project a manager
oversees, while Pulse's per-project cost stays close to zero, so the benefit
increases with portfolio size. A second benefit, earlier detection of a
deteriorating project, is deliberately not claimed in monetary terms: the system
has not run long enough to observe an instance, and claiming it would be
unsupported.

\paragraph{Revenue model.} The comparable commercial products identified in
\Cref{ch:research} --- LinearB, Swarmia, Faros AI and Jellyfish --- are priced
per seat. Per-seat pricing sits poorly with Pulse's cost structure for two
reasons. First, cost scales with tracked projects and survey cycles rather than
with logins, so per-seat pricing would systematically misprice the product
against its own cost base. Second, and more seriously, charging per developer
creates a direct financial incentive to exclude developers from the surveys, and
the representativeness of that survey signal is the product's principal
differentiator. A \textbf{per-tracked-project subscription} aligns price with
cost and leaves survey coverage free at the margin, and is therefore the model
proposed.

\subsection{Sustainability in Business and Commercialisation}
\label{subsec:commercialisation}

The following points determine the system's future after the project.

\begin{itemize}
  \item \textbf{Handover.} The system is hosted and operational (C-08) and will
  be handed over to SSCL.
  \item \textbf{Custody of infrastructure.} The API and worker run on a
  free-tier Render deployment (\Cref{subsec:deviation-deployment}), with
  Supabase, Redis, Vercel, Gemini, Pinecone and Brevo accounts held by the team.
  At handover each of these will be transferred to or recreated under SSCL's
  ownership, and every credential rotated --- including the AES key protecting
  per-project connector credentials, since re-keying requires re-encrypting the
  stored values rather than simply changing an environment variable. The
  encrypted connector credentials in the database are SSCL's property, not the
  team's.
  \item \textbf{NDA obligations.} The work is governed by a non-disclosure
  agreement signed with SSCL (C-09). Because evaluation used the team's own
  repositories, no SSCL production data is held in the team's database.
  \item \textbf{What production would require beyond the current deployment.} A
  paid database tier with backups; a redundant or restart-supervised worker,
  since a single worker instance is a single point of failure for all
  synchronisation, survey dispatch and embedding; centralised log aggregation
  and alerting (gap G-04); and a named on-call owner. The free-tier deployment
  is adequate to demonstrate the system but not to run it for a customer at
  scale.
\end{itemize}

\section{Individual and Teamwork}
\label{sec:teamwork}

\subsection{Participation and Contribution}
\label{subsec:contribution}

The team comprised five students. \Cref{tab:contributions} records each member's
role and the modules they worked on.

\begin{longtable}{
  >{\raggedright\arraybackslash}p{0.20\textwidth}
  >{\raggedright\arraybackslash}p{0.13\textwidth}
  >{\raggedright\arraybackslash}p{0.18\textwidth}
  >{\raggedright\arraybackslash}p{0.34\textwidth}
}
\caption{Team roles and modules}
\label{tab:contributions}\\
\toprule
\textbf{Name} & \textbf{Student ID} & \textbf{Role} & \textbf{Modules worked on} \\
\midrule
\endfirsthead
\toprule
\textbf{Name} & \textbf{Student ID} & \textbf{Role} & \textbf{Modules worked on} \\
\midrule
\endhead
\bottomrule
\endlastfoot
\guidance{[Name]} & \guidance{[ID]} & \guidance{[Role]} & \guidance{[Modules]} \\
\addlinespace
\guidance{[Name]} & \guidance{[ID]} & \guidance{[Role]} & \guidance{[Modules]} \\
\addlinespace
\guidance{[Name]} & \guidance{[ID]} & \guidance{[Role]} & \guidance{[Modules]} \\
\addlinespace
\guidance{[Name]} & \guidance{[ID]} & \guidance{[Role]} & \guidance{[Modules]} \\
\addlinespace
\guidance{[Name]} & \guidance{[ID]} & \guidance{[Role]} & \guidance{[Modules]} \\
\end{longtable}

\subsection{Collaboration and Conflict Resolution}
\label{subsec:collaboration}

Collaboration was structured around a weekly meeting. At each meeting, tasks
were assigned for the coming week. At the next meeting, each member reported how
far their assigned work had progressed, and the team checked for dependencies
between tasks --- whether one member's work was waiting on another's --- so that
blocked work was identified and re-sequenced early rather than discovered at
integration. Technical disagreements were raised and settled in the same
meeting; where a question affected the product's direction, it was taken to the
industry collaborator (\Cref{subsec:leadership}).

Changes reached \texttt{main} through pull requests. Because C-02 left the team
without dedicated quality-assurance or operations members, review was the point
at which a second person saw a change before it reached \texttt{main}, and the
automated gate of NFR-23 existed partly so that review effort could be spent on
design and intent rather than on correctness a machine can establish.

\subsection{Leadership and Direction}
\label{subsec:leadership}

Technical direction was set through a cycle of internal discussion and partner
review. The team discussed possible improvements among themselves, implemented
them, and then presented each major improvement to SSCL's Chief Technical
Officer, who gave technical direction for the next stage of work
(\Cref{subsec:multidisciplinary}). The team therefore owned the design of each
change, while the partner set its priorities and kept the product aligned with
what its intended users needed.

\Cref{ch:evaluation} supplies a ready example of a setback converted into a
revised plan, and it is stronger evidence of direction than any general
statement about leadership. When transactional email proved unreliable, the
response was to migrate the provider and eliminate the failure mode
(\Cref{subsec:deviation-email}) rather than to add retries around a mechanism
that was structurally unsuited to the job.

\subsection{Multidisciplinary Engagement}
\label{subsec:multidisciplinary}

The project required the team to work outside software engineering proper in
three distinct directions.

\paragraph{With the industry partner.} Requirements came from SSCL's Chief
Technical Officer and from an intern with operational exposure, which meant
translating a managerial complaint about dashboard reconciliation into a
specification with verifiable acceptance criteria --- a requirements-engineering
and client-management activity rather than a programming one. Engagement
continued throughout development: the team held five meetings with SSCL's CTO,
each after a major improvement to the system, and the design changed as a
result:

\begin{itemize}
  \item \textbf{Audience.} The system was initially conceived for developers as
  well as leadership. SSCL advised designing it only for top-level users such as
  the CTO or CEO, which fixed the product's positioning as a portfolio-level
  management tool.
  \item \textbf{Anonymity.} SSCL required surveys to be anonymous, which the
  design enforces in the data model (NFR-07, \Cref{subsec:inv2}).
  \item \textbf{Delivery.} SSCL asked that survey links be sent by email rather
  than posted to Slack, Discord or any other public channel. Email delivery to
  each eligible developer is the default; chat-channel broadcast is used only if
  an organisation chooses to supply a Slack, Discord or Telegram token.
  \item \textbf{Frequency.} SSCL asked that a developer receive no more than one
  survey within a set period --- a month, or 15 days --- and that the period be
  configurable. This is the cooldown of FR-33
  (\texttt{SURVEY\_MIN\_DAYS\_BETWEEN\_SURVEYS}, default 15 days).
  \item \textbf{Action logging.} The team proposed the action log, and SSCL
  approved it; it became FR-37--FR-41.
\end{itemize}

\paragraph{Into the measurement of work itself.} Deciding what constitutes
``engineering health'', and how team morale and blockers can be measured without
measuring individuals, draws on organisational behaviour and survey methodology
rather than on computing. Several of the project's most consequential decisions
are of this kind and not technical at all: collecting qualitative signal
anonymously; declining to let a manager steer question wording
(\Cref{subsec:deviation-guidance}); and designing the effectiveness review so
that an owner may defer judgement rather than being forced to rate an outcome
before evidence exists. Each trades data quality or convenience for measurement
validity or respondent welfare.

\paragraph{Into research ethics and data protection.} Deciding what may
legitimately be inferred about identifiable people from data they did not
volunteer, and what may be transmitted to a third-party model, is a governance
question. The GDPR-derived principles treated as design constraints in
\Cref{ch:requirements}, and the adversarial anonymity analysis in
\Cref{subsec:inv2}, are the output of that work.

\section{Impact on Society}
\label{sec:society}

The system's users are engineering managers, but the people it measures are
developers, and the two groups experience it very differently. This section
considers both, and covers health, safety, cultural and legal issues explicitly.

\paragraph{Benefit: earlier detection.} The primary benefit is that a
deteriorating project is identified from accumulated trend data rather than from
its first visible failure --- a missed deadline or a broken release. Those
affected are the client's leadership and managers (SH-01, SH-02) and,
indirectly, the developers, for whom an intervention arriving before a crisis is
less disruptive than one arriving after.

\paragraph{Benefit: a feedback channel that did not previously exist.}
Developers (SH-03) gain a route by which low morale, persistent blockers or loss
of confidence in a plan reach management without any individual having to raise
it personally. \Cref{ch:problem} identifies the absence of this channel as one
of the two gaps the project addresses, and the anonymity of the channel is
precisely what makes using it safe. This benefit is therefore \textbf{contingent
on the anonymity guarantee being real}, which is why anonymity is enforced in
the data model itself rather than promised in policy (\Cref{subsec:inv2}).

\paragraph{Health and safety.} The system carries no physical safety
implication: it controls no equipment and operates strictly read-only against
the partner's tools (C-07). Its health dimension is occupational and
psychological. Here the design has one genuine risk to answer for, addressed
below, and one deliberate protection: the effectiveness review (FR-40) permits
an action owner to defer judgement rather than forcing a rating before evidence
exists, and presents overdue items in deliberately non-alarming styling ---
amber rather than red --- with a persistent inline banner rather than a modal
that interrupts or a toast that vanishes on a timer. Both choices accept a
weaker dataset in exchange for not manufacturing pressure, which is the correct
trade when the pressure would fall on a person and the data is only diagnostic.

\paragraph{Principal risk: repurposing as individual surveillance.} A system
that scores team health and reports on engineering practice can be turned
against the people it measures. Three misuses are foreseeable:

\begin{enumerate}
  \item \emph{Reading a low team-health score as an indictment of named
  individuals.} This cannot be prevented technically --- a manager determined to
  read a project score as a verdict on a person will do so. The available
  mitigation is presentational: because every score expands into the metrics
  that produced it (FR-26), a manager who inspects a declining score sees bus
  factor, ownership concentration or review network density, not a name.
  \item \emph{Using survey findings to identify and pressure dissenters.}
  Answered structurally by omitting respondent identity from
  \texttt{survey\_response}: no stored response can be traced to its author.
  Statistical inference from timing or content remains possible in principle
  (\Cref{subsec:inv2}).
  \item \emph{Treating connector metrics such as commit or review counts as
  individual productivity measures.} Answered by scoring only at project level
  and by excluding per-individual measurement from scope explicitly in
  \Cref{ch:problem}. Note that the raw data to do this nonetheless passes
  through the system, so the protection is a design commitment rather than an
  impossibility; a future maintainer could add per-developer views without
  fighting the architecture.
\end{enumerate}

\paragraph{Cultural considerations.} Survey candour depends on trust that
anonymity is real, and that trust is cultural at least as much as technical. In
a hierarchical workplace, or a small team where any response narrows authorship,
developers may reasonably decline to answer candidly whatever the schema
guarantees. This bounds the system's usefulness in exactly the situations where
team health is most at risk. Enforcing anonymity in the schema is the technical
half of that trust; the organisational half --- management visibly not
attempting to identify respondents --- lies with the adopting organisation.

\paragraph{Legal considerations.} Survey responses relate to identifiable
individuals' working conditions even where the responses themselves are
unattributable, and connector metrics include developer identities in their raw
form. The design response is documented in \Cref{ch:requirements}: GDPR
principles treated as design constraints, prompts constructed to exclude
personal names and work-item identifiers, and credentials encrypted at rest.
This evaluation verified the prompt claim by inspection --- no author, assignee,
login, email, issue-key or branch fields appear in the AI payload construction
(\Cref{subsec:eval-nfr}, NFR-09). The residual legal exposures are the free-text
answers transmitted verbatim (gap G-03) and the absence of any data-subject
access, rectification or erasure mechanism (gap G-01).

\section{Environmental Impact and Life Cycle Sustainability}
\label{sec:environment}
\label{sec:sustainability}

The system owns no hardware, so its footprint is entirely the third-party cloud
capacity it consumes, and it scales with usage rather than being fixed. Impact
is considered across three life-cycle stages, with mitigation proposed for each
negative impact.

\paragraph{Development.} Consumption comprised developer workstations and the CI
pipeline. CI is the non-obvious contributor: it runs type checking, linting, a
build and two test suites on every push, and across a year of development that
accumulates. Two design decisions already reduce it. The path filter means a
backend-only change never builds the frontend, and the \texttt{concurrency}
group with \texttt{cancel-in-progress} means that pushing twice in quick
succession cancels the superseded run rather than letting both complete. Both
were adopted for speed and reduce compute proportionally as a side effect.

\paragraph{Operation.} Three components dominate. Scheduled synchronisation runs
per project per configured slot and issues bounded-concurrency API calls; AI
inference is the most energy-intensive operation per unit, covering embeddings
for action search and two generation calls per survey cycle; and the always-on
worker consumes continuously regardless of load, which at small portfolio sizes
makes it the largest single share of an otherwise small total.

Several existing design decisions reduce operational impact, and each is worth
naming because each was taken for a different reason and has this as a side
effect:

\begin{itemize}
  \item Embeddings are computed once per action and keyed by
  \texttt{content\_hash}, so an unedited action is never re-embedded and a
  repeated search never recomputes an embedding the system already holds.
  \item Ordinary lexical search removes AI inference from the request path when
  the user selects that mode. It remains available during an AI outage, although
  a deep-search request currently returns 503 instead of switching to it
  automatically (\Cref{subsec:inv3}).
  \item Dashboard reads are served from stored snapshots and never trigger live
  connector calls (NFR-17), so viewing a dashboard, the most frequent user
  action, costs nothing externally.
  \item Synchronisation is scheduled rather than continuous, which follows
  directly from the product's positioning as a trend indicator and is the single
  largest avoided cost relative to a real-time design.
\end{itemize}

Two further reductions are available and are proposed here even though they are
not implemented: \textbf{synchronising only projects whose upstream state has
changed} since the last snapshot, detectable cheaply from a repository's latest
commit SHA before issuing the full metric fetch; and \textbf{pruning or
aggregating snapshot detail beyond a retention horizon}, since a two-year-old
daily snapshot serves no analytical purpose that a weekly aggregate would not,
and storage grows unboundedly at one snapshot per project per day.

\paragraph{Retirement.} Two distinct retirements must be handled and neither
currently has an implemented path.

When a \emph{customer offboards}, their connector credentials, snapshots,
scores, surveys, responses and action log should be deleted on a stated
timetable, and their access tokens revoked at the provider so that they cannot
be used even if a copy of the ciphertext survives. When the \emph{system itself}
is retired, the same applies across all tenants, together with deletion of
embeddings held in the external Pinecone index --- which lives outside the
primary database and is easily overlooked precisely because deleting the
Postgres rows feels like completing the job.

\textbf{Neither path exists today.} Gap G-02 records that no retention schedule
is defined, and G-01 that no erasure mechanism is exposed. This is
simultaneously a sustainability gap and a data-protection exposure, and it is
required work before any handover to SSCL.

\section{Ethics}
\label{sec:ethics}

\subsection{Equity and Inclusivity}
\label{subsec:equity}

\paragraph{Accessibility.} The frontend is built on Radix UI primitives, which
supply keyboard interaction and ARIA semantics for dialogs, menus, tabs and
similar components by default, so a baseline is inherited rather than earned.
Beyond that baseline, this evaluation found deliberate accessibility work in the
codebase: 46 \texttt{aria-label} attributes, 23 explicit \texttt{role}
assignments, and uses of \texttt{aria-hidden}, \texttt{aria-expanded},
\texttt{aria-modal}, \texttt{aria-labelledby}, \texttt{aria-haspopup},
\texttt{aria-checked} and \texttt{aria-valuenow}, plus one visually-hidden
(\texttt{sr-only}) heading. That is more than an untouched component library
provides and indicates the concern was actively considered.

What has \textbf{not} been established is conformance. No automated audit or
assistive-technology testing was performed; colour contrast ratios have not been
measured against the WCAG Level AA threshold; and the Recharts trend
visualisations offer limited screen-reader semantics. A specific risk worth
checking is that risk level is bucketed into LOW/MEDIUM/HIGH at the 40/70
thresholds and rendered with colour coding --- a red/amber/green scale is the
classic failure mode for colour-vision deficiency unless the numeric value is
always rendered alongside it. Conformance is therefore claimed as
\textbf{partial and unverified}, consistent with gap G-06, and attributable to
C-02: the team had no dedicated design or accessibility specialist.

\paragraph{Fairness across teams of different shapes.} The system must not
advantage teams that happen to use the tools it integrates best. Two provisions
address this. Weight renormalisation (FR-22) means a team without SonarQube or
without Jira is scored on what it does have rather than penalised for what it
lacks --- and \Cref{subsec:inv1} shows what the alternative would have cost such
a team: on the worked example, a security score of 12 rather than 80, decided
entirely by which tools were connected rather than by how the team works. The
connector abstraction (NFR-21) means a team on a different toolchain is excluded
by implementation effort rather than by design, as the six shipped connectors
demonstrate.

\paragraph{Who is excluded.} Two groups, stated plainly. Teams whose tools have
no connector cannot be scored at all. And newly adopted projects have no trend
data at all, since C-06 records that historical backfill is impossible, so the
system is least useful precisely when a team first tries it and must be
persuaded of its value.

\subsection{Accountability and Personal Responsibility}
\label{subsec:accountability}

\paragraph{For the system's output.} Pulse produces diagnostic signals about
projects, not judgements about people, and that distinction is load-bearing
rather than rhetorical. A score identifies where to look; accountability for
what is done with that information rests with the manager who acts on it, not
with the tool. The design supports this rather than merely asserting it: every
score is decomposable into the metrics that produced it (FR-26, NFR-20), so a
manager cannot honestly treat a number as an unexaminable verdict, and the
deliberate separation of the connector-derived risk score from the AI-derived
survey insight means neither silently contaminates the other's meaning.

\paragraph{Within the team.} Every change reached \texttt{main} through a
reviewed pull request, so each change carries both an identified author and an
identified reviewer, and the automated gate of NFR-23 meant neither could merge
work that failed to build or to pass its tests. Module responsibility followed
task assignment at the weekly meeting (\Cref{subsec:collaboration}).

\paragraph{A note on this evaluation.} The anonymity analysis in
\Cref{subsec:inv2} was conducted by the team adversarially, against its own
system. Accountability for a system's behaviour includes testing its guarantees
from the attacker's side, particularly where the guarantee is a promise made to
third parties --- in this case SSCL's developers, whose survey responses are
promised anonymity.

\subsection{Proper Use of Intellectual Property}
\label{subsec:ip}

\paragraph{Third-party software.} The system is built on open-source
dependencies consumed through npm. On the backend: Express,
\texttt{@supabase/supabase-js}, BullMQ, ioredis, Octokit (with the retry and
throttling plugins), \texttt{@slack/web-api}, \texttt{@google/genai}, bcryptjs,
helmet, cors, cookie-parser, morgan, pino, express-rate-limit, fast-xml-parser,
adm-zip, nodemailer and resend. On the frontend: React, Vite, React Router,
Tailwind CSS, Radix UI, MUI (\texttt{@mui/material}, \texttt{@mui/icons-material}),
Recharts, Emotion, \texttt{react-hook-form}, \texttt{date-fns},
\texttt{lucide-react}, \texttt{next-themes}, \texttt{motion}, \texttt{cmdk},
\texttt{sonner}, \texttt{react-dnd}, \texttt{canvas-confetti} and others, with
ESLint, TypeScript and Vitest in the toolchain.

\paragraph{Third-party services.} GitHub, GitLab, Jira, Linear, SonarCloud,
Gemini, Pinecone, Brevo, Slack, Telegram and Discord are used through their
public APIs under their respective terms of service, read-only against the
partner's instances using tokens the partner supplies (C-07). SH-06 records
adherence to those terms as a stakeholder expectation, and the rate-limit
provisions verified in \Cref{subsec:inv4} are part of honouring it rather than
merely of staying within a budget.

\paragraph{Partner material.} Work is conducted under a non-disclosure agreement
signed with SSCL (C-09) governing both the handling of their data and what this
report may disclose.

\paragraph{Generative AI.} The candidates' declaration requires disclosure of
generative-AI use, and it should be stated specifically and completely,
separating the \emph{product} from the \emph{process}.

\begin{itemize}
  \item \textbf{In the product:} Google Gemini is a runtime component of the
  delivered system. It generates survey questions, scores survey responses and
  generates the insight summary, and asynchronously produces stored action
  embeddings (\texttt{gemini-embedding-001}). The current deep-search request
  path reranks action candidates directly through Pinecone; INV-3 found that it
  does not fall back when Pinecone is unavailable. This is disclosed as an
  architectural dependency, not as authoring assistance, and is described
  throughout \Cref{ch:design,ch:evaluation}.
  \item \textbf{In the process:} AI coding assistants, principally Claude and
  Codex, were used for coding and debugging. In every case the team reviewed and
  verified the output and takes responsibility for it.
\end{itemize}

\subsection{Professionalism and Ethical Codes}
\label{subsec:codes}

Three principles of the \textbf{ACM Code of Ethics and Professional Conduct}
bore directly on decisions recorded elsewhere in this report. They are given
here with the decision each governed and the cheaper alternative that was
rejected, because that is what distinguishes engagement with a code from
citation of one.

\paragraph{Respect privacy (ACM 1.6).} The Code requires that systems collecting
personal information collect only what is necessary and not repurpose it beyond
what the subject understood. Pulse collects opinions about a project's condition
from people whose employer will read the result. The decision this principle
governed was \emph{where} to enforce anonymity. Omitting identity from
\texttt{survey\_response} means direct attribution would require a deliberate
schema change rather than an accidental application join. Investigation INV-2
then tested that protection adversarially and found every stored linkage route
closed. ACM 1.3 (honesty about limitations) is why the statistical routes that
no schema can close are reported alongside that result rather than omitted.

\paragraph{Avoid harm (ACM 1.2) and evaluate the risks of systems (ACM 2.5).}
The foreseeable harm from this system is not technical failure but correct
operation misapplied: scores read as verdicts on individuals, survey findings
used to identify dissenters. ACM 2.5 requires comprehensive risk evaluation
rather than a favourable one, which is why \Cref{sec:society} records the misuse
the design \emph{cannot} prevent alongside the two it can, and why
\Cref{subsec:inv2} was conducted as an adversarial exercise against the team's
own system rather than as a confirmation of it.

\paragraph{Hold the public good as the central concern (ACM 3.1; IEEE Code of
Ethics clause 1).} This principle settled the two decisions where user welfare
and engineering convenience pulled in opposite directions. Removing the
manual-guidance override (\Cref{subsec:deviation-guidance}) discarded working,
client-requested functionality because it would have quietly biased what the
survey measured, and therefore what developers were understood to have said. The
deferrable effectiveness review (FR-40) accepts a weaker dataset rather than
pressing an owner to certify an outcome before the evidence exists. In both
cases the convenient option was available, defensible on delivery grounds, and
not taken.

\paragraph{Institution of Engineers, Bangladesh (IEB) code of ethics.} The IEB
code requires an engineer to act with integrity towards clients and employers
and to respect the confidentiality of the information they entrust. That
obligation governed how the project treated its industry partner: the work was
carried out under a signed non-disclosure agreement, and the system was
evaluated against the team's own repositories rather than SSCL's. Testing on
partner data would have been quicker and more realistic, but it would have
brought SSCL's confidential engineering data into a student development
environment, and it was not done.
